% Extrait : les points sommaires de la version française du
% sommaire sur la demande d'accès 2026-581, mot pour mot. pdflatex (deux fois).
\documentclass[10pt,letterpaper]{article}
\usepackage[T1]{fontenc}
\usepackage[utf8]{inputenc}
\usepackage[scaled=0.95]{helvet}
\renewcommand{\familydefault}{\sfdefault}
\usepackage[letterpaper,top=0.8in,bottom=0.8in,left=0.9in,right=0.9in]{geometry}
\usepackage{microtype}
\usepackage{enumitem}
\usepackage{titlesec}
\usepackage{fancyhdr}
\usepackage{xcolor}
\usepackage{booktabs}
\usepackage{longtable}
\usepackage{array}
\usepackage[hidelinks]{hyperref}
\urlstyle{same}
\newcommand{\docversion}{2026-10-11-14-15-00} % set by build.py (America/Toronto)
\hypersetup{pdfsubject={version \docversion}}

\setlength{\parindent}{0pt}
\setlength{\parskip}{0.5em}
\linespread{1.0}

\titleformat{\section}{\normalfont\large\bfseries}{\thesection}{0.6em}{}[\vspace{1pt}\titlerule]
\titlespacing*{\section}{0pt}{1.1em}{0.4em}
\titleformat{\subsection}{\normalfont\bfseries}{\thesubsection}{0.6em}{}
\titlespacing*{\subsection}{0pt}{0.8em}{0.25em}

\pagestyle{fancy}
\fancyhf{}
\renewcommand{\headrulewidth}{0pt}
\fancyfoot[L]{\footnotesize\color{gray}Demande d'accès 2026-581 -- points sommaires -- \pagesurl}
\fancyfoot[R]{\tiny\color{gray}v.\,\docversion}

% ---- Repository links: fill in once the documents are uploaded ----
\newcommand{\repourl}{https://github.com/OpenMobilityData/Terrebonne-access-requests}
\newcommand{\repo}[1]{\href{\repourl}{#1}}
% \doc{path-in-repo}{label}: link to a specific file in the repository
\newcommand{\doc}[2]{\href{\repourl/blob/main/#1}{#2}}
% GitHub Pages URL (Settings > Pages > deploy from main, root). Page anchors
% (#page=N) work there because Pages serves PDFs with the right content type;
% GitHub's blob viewer and raw.githubusercontent.com do not honour them.
\newcommand{\pagesurl}{https://openmobilitydata.github.io/Terrebonne-access-requests}
\expandafter\def\csname bundle01\endcsname{2026-581/raw/01_emails_mandate_counts.pdf}
\expandafter\def\csname bundle02\endcsname{2026-581/raw/02_framework_agreement_tender.pdf}
% \pg{bundle}{pages}: page reference, linked to the bundle at its first page
% \pgl{bundle}{anchor page}{displayed pages}: same with an explicit anchor
\newcommand{\pgl}[3]{\href{\pagesurl/\csname bundle#1\endcsname\#page=#2}{\textsf{\small[#1, p.~#3]}}}
\newcommand{\pg}[2]{\pgl{#1}{#2}{#2}}
\newcommand{\fillin}[1]{\textcolor{red}{[#1]}}
\newcommand{\quot}[1]{\begin{quote}\small #1\end{quote}}

\begin{document}

{\large\bfseries Demande d'accès 2026-581~: ce que révèle le dossier de l'arrondissement}\hfill{\small 9~octobre 2026, rév.\ 11~octobre}

\vspace{-0.2em}
{\small Points sommaires du sommaire et analyse, mot pour mot. Analyse complète, références aux pages et documents communiqués~:\\ \repo{\repourl}}
\vspace{0.4em}

\begin{enumerate}[nosep,itemsep=0.3em,leftmargin=1.4em]
  \item \begin{minipage}[t]{\linewidth}\textbf{La décision a précédé l'étude.} Le 10~juin 2026, la
    division de l'ingénierie de l'arrondissement écrivait à la division de
    l'aménagement des rues de la Ville que les élus souhaitaient revoir
    l'aménagement de 2024 «\,avec l'objectif de redonner la bande de
    stationnement qui avait été retranchée\,», que les ingénieurs de
    l'arrondissement avaient \emph{déjà} émis un avis technique sur les deux
    concepts esquissés par les élus, et que ceux-ci voulaient «\,un avis
    officiel, neutre et externe\,». Le mandat finalement confié à EXP le
    20~août a pour titre, en partie, «\,évaluer la place du stationnement sur
    rue, dans le but d'en augmenter le nombre\,».\end{minipage}\vspace{0.35em}
\item \begin{minipage}[t]{\linewidth}\textbf{L'équipe vélo de la Ville elle-même a soulevé des enjeux
    de sécurité.} Le Service de l'urbanisme et de la mobilité a remis une
    analyse en mode express le 17~juin\,; le résumé qu'en fait
    l'arrondissement, rédigé cinq jours plus tard, indique que les options
    «\,présentent plusieurs compromis importants sur le plan de la sécurité,
    du confort des usagers vulnérables et de la cohérence du réseau
    cyclable\,». Cette analyse n'a pas été communiquée. L'avis du Service de
    sécurité incendie du 13~juillet est caviardé en entier.\end{minipage}\vspace{0.35em}
\item \begin{minipage}[t]{\linewidth}\textbf{La mesure hivernale est qualifiée de transitoire\,; le
    mandat du cabinet de la mairesse était permanent.} Le 13~août, la cheffe
    de cabinet de la mairesse d'arrondissement a fixé le mandat par écrit~:
    une piste cyclable bidirectionnelle d'un côté de Terrebonne, de Loyola à
    Côte-Saint-Luc, «\,qui respecte les normes minimales prévues par la loi\,»,
    et «\,le retour du stationnement d'un côté de la rue\,». Aucune mesure
    hivernale n'y figure. La demande de services, une semaine plus tard, a
    ajouté un Volet~1, le retrait des aménagements cyclables en hiver, à
    titre de «\,mesure transitoire pour l'hiver 2026--2027\,», avant le
    Volet~2, la piste bidirectionnelle permanente. La conversion hivernale
    annoncée le 6~octobre est la première étape d'une réduction à l'année de
    l'espace cyclable qui n'a jamais été soumise au public.\end{minipage}\vspace{0.35em}
\item \begin{minipage}[t]{\linewidth}\textbf{L'«\,étude\,» est une note technique de six semaines et de
    14\,755~\$, avec une date de conseil fixée d'avance.} Demande de
    services le 20~août\,; offre le 25~août\,; bon de commande le 26~août\,;
    réunion de démarrage le 3~septembre\,; version préliminaire le
    17~septembre\,; version finale due le 2~octobre\,; annonce le 6~octobre. Le
    personnel de l'arrondissement a écrit au consultant le 18~août que
    «\,l'objectif de nos élus est de pouvoir agir sur vos conclusions pour un
    horizon début octobre\,» et, le 10~septembre, que la version préliminaire
    devait être livrée au plus tard le 18 «\,pour respecter notre calendrier
    du Conseil d'arrondissement du 5~octobre\,».\end{minipage}\vspace{0.35em}
\item \begin{minipage}[t]{\linewidth}\textbf{Les comptages de vélos cités au conseil reposent sur les
    données les plus faibles du dossier et sous-estiment l'usage
    hivernal\,; le dossier contenait les données du compteur permanent pour
    l'hiver dernier, qui attestent un usage hivernal moyen au moins cinq fois
    plus élevé que le chiffre cité. La présentation n'a pas tenu compte de
    ces données plus robustes.} Le «\,12\,» et le «\,8\,» proviennent de la
    présentation des comptages de l'arrondissement de février 2025~: il ne
    s'agit ni de vélos par jour ni de vélos par heure, mais de vélos comptés
    au moyen d'une caméra temporaire (Miovision) installée aux intersections
    Cavendish et Grand pendant six heures de pointe d'une seule journée, le
    29~janvier 2025, lors du premier hiver de la piste, alors que le
    déneigement était souvent en retard, et lus sur une diapositive qui ne
    précise rien de tout cela. Au moment de la séance, un compteur permanent
    avait enregistré chaque vélo passé sur la rue pendant l'hiver suivant,
    soit une année entière plus récente\,; les ingénieurs de l'arrondissement
    avaient offert ses données à la Ville en juin, et son tableau était entre
    les mains du consultant en août~: 41 par jour de décembre à mars (46 les
    jours de semaine), et un deuxième été à environ 700 par jour. La baisse
    de «\,99~\%\,» citée à la séance, 700 contre 8, oppose le plus grand nombre
    d'une diapositive, Grand en septembre 2024 compté sur les quatre
    approches, au plus petit nombre de la suivante, Grand en janvier 2025 sur
    l'axe Terrebonne seulement\,; comparées de façon équivalente, les
    diapositives elles-mêmes donnent 97 à 98~\%, étiquettent janvier à
    $-95$~\%, et le compteur donne environ 93~\% d'une saison entière à
    l'autre. Les deux sources s'entendent sur le fait que la baisse est
    forte\,; elles diffèrent sur le nombre absolu, et une seule d'entre elles
    donne «\,quatre personnes\,». Quelle qu'en soit l'explication, une décision
    présentée comme fondée sur les données reposait sur six heures
    de comptage en une seule journée et a ignoré des milliers d'heures de
    comptage plus récent sur tout un hiver.\end{minipage}\vspace{0.35em}
\item \begin{minipage}[t]{\linewidth}\textbf{Le côté stationnement du bilan n'a jamais été mesuré.}
    Dans les 508~pages, il n'y a aucun inventaire des cases sur rue, aucun
    relevé d'occupation, aucun sondage auprès des résidents et aucune
    estimation du nombre de cases que la conversion créerait\,; les «\,50 à
    60 cases\,» que le conseiller a données à l'APC-NDG lors de la rencontre
    de consultation ne figurent nulle part dans les documents
    communiqués. Le mandat demande au consultant de compter les cases qu'un
    aménagement «\,pourrait permettre de recréer\,» et de mesurer
    «\,l'amélioration de la disponibilité du stationnement\,»\,; l'offre
    promet une analyse de «\,l'impact sur le nombre de places\,» et ne
    prévoit aucun travail de terrain\,; le consultant a travaillé à partir
    d'imagerie aérienne. Ce que la note dit du stationnement est donc un
    décompte de cases qu'on pourrait tracer, non une mesure du besoin. Rien
    dans le dossier n'indique combien de nuits d'hiver le stationnement
    près de Terrebonne est réellement complet.\end{minipage}\vspace{0.35em}
\item \begin{minipage}[t]{\linewidth}\textbf{L'argument en faveur du rétrécissement repose sur l'idée
    que les pistes sont plus larges que leur circulation cycliste ne le
    justifie\,; cette idée est l'opinion des élus, non l'avis des
    ingénieurs.} Le courriel du 10~juin rapporte la préoccupation des élus
    selon laquelle la largeur des pistes était «\,jugée disproportionnée par
    rapport au volume de cyclistes observé \ldots\ selon les élus\,», et le
    conseiller Teodoresco a repris le point devant l'APC-NDG lors de la
    rencontre de consultation, mais non à la séance du conseil du
    6~octobre. Aucun ingénieur au dossier ne l'endosse. Les ingénieurs de
    l'arrondissement décrivent les solutions de rechange comme «\,une voie de
    circulation très étroite (moins de 4~mètres)\,» et «\,un aménagement
    bidirectionnel très étroit\,»\,; l'équipe vélo de la Ville y a vu
    «\,plusieurs compromis importants\,»\,; et le mandat du 13~août demande
    un aménagement «\,qui respecte les normes minimales prévues par la
    loi\,», ce qui revient à admettre que les concepts de juin ne les
    respectaient pas. La largeur d'une piste est fixée par le gabarit du
    cycliste, l'espace de dépassement et le matériel de déneigement, dont
    aucun ne dépend du nombre de personnes qui roulent\,; rien au dossier ne
    décrit les pistes actuelles comme plus larges que la norme ne
    l'exige.\end{minipage}\vspace{0.35em}
\item \begin{minipage}[t]{\linewidth}\textbf{Une option limitée aux opérations de déneigement
    faisait partie du mandat.} L'offre d'EXP énumère, parmi les scénarios à
    comparer, «\,la fermeture en période de déneigement seulement\,». Le
    conseiller a affirmé à la séance du 6~octobre que le stationnement limité
    aux tempêtes avait été écarté\,; l'analyse sur laquelle repose cette
    affirmation se trouve dans la note qui n'a pas été
    communiquée.\end{minipage}\vspace{0.35em}
\item \begin{minipage}[t]{\linewidth}\textbf{L'exercice ne suit pas la méthode que la Ville a elle-même
    publiée pour réviser une piste cyclable.} Le 10~février 2026,
    l'administration a lancé son audit du réseau cyclable sur le principe
    «\,recueillir les données, analyser et corriger au besoin\,»~: analyse
    technique, collecte de données et observation terrain pendant la haute
    saison, sous la direction du Service de l'urbanisme et de la mobilité,
    avec des résultats attendus à la fin de 2026\,; le même communiqué nomme
    Terrebonne parmi les rues où la Ville travaille avec l'arrondissement sur
    la sécurité. Sur Terrebonne, l'objectif a été fixé avant toute analyse,
    l'analyse du SUM a été défavorable et refusée, le mandat visait à
    augmenter le stationnement, aucune donnée n'a été recueillie, et le
    changement a été annoncé alors que l'audit de la Ville était toujours en
    cours.\end{minipage}\vspace{0.35em}
\item \begin{minipage}[t]{\linewidth}\textbf{Le refus de communiquer le rapport technique repose sur un
    motif qui s'éteint avec la décision.} La note et l'analyse de la Ville sont
    refusées en vertu de l'article~39, qui s'applique «\,jusqu'à ce que la
    recommandation ait fait l'objet d'une décision\,». Le conseiller a annoncé
    la décision le 6~octobre et promis des changements d'ici Noël. Ou bien la
    décision a été prise et le motif n'existe plus, ou bien elle ne l'a pas
    été et l'annonce était prématurée. Une demande de révision peut être
    adressée à la Commission d'accès à l'information dans les 30~jours, soit
    d'ici le \textbf{8~novembre 2026}.\end{minipage}\vspace{0.35em}
\end{enumerate}

\end{document}
